\chapter{Kết luận và Hướng phát triển}
\label{chap:conclusion}

\section{Kết luận}
\label{sec:conclusion}

Khóa luận đã nghiên cứu và xây dựng thành công \textbf{ViNeoBERT} -- một mô hình
Encoder hiện đại cho tiếng Việt -- bằng cách thích nghi NeoBERT thông qua kỹ thuật
SALT được cải tiến, thay vì huấn luyện lại từ đầu. Đối chiếu với các mục tiêu đặt ra
ở Chương~\ref{chap:intro}, đề tài đã đạt được các kết quả sau:

\begin{itemize}
  \item \textbf{Về nghiên cứu kiến trúc}: Chúng tôi đã phân tích các cải tiến kiến
        trúc của NeoBERT (RoPE, SwiGLU, Pre-RMSNorm, tỷ lệ sâu--hẹp tối ưu) và xác
        định đúng hai điểm mấu chốt ảnh hưởng tới việc chuyển giao: chuẩn hóa
        Pre-RMSNorm nhạy với độ lớn của nhúng -- dẫn tới bước hiệu chỉnh chuẩn -- và
        đầu giải mã tách rời đòi hỏi khởi tạo riêng cho lớp đầu ra
        (Chương~\ref{chap:related}, \ref{chap:method}).

  \item \textbf{Về phương pháp chuyển giao}: Chúng tôi đã triển khai thuật toán SALT
        và cải tiến nó cho cặp ngôn ngữ Anh--Việt với ba đóng góp: (1)~mở rộng tập neo
        bằng cặp dịch khai thác tự động, kèm cơ chế lọc từ đồng tự khác nghĩa đặc thù
        cho ngôn ngữ hệ chữ Latin; (2)~áp dụng phép chiếu theo từng token cho cả lớp
        đầu ra tách rời, kết hợp khởi tạo độ chệch theo tần suất từ đơn và hệ số co
        trọng số $\gamma = 0{,}1$; (3)~bổ sung giai đoạn căn chỉnh đóng băng trước khi
        tiếp tục tiền huấn luyện (Chương~\ref{chap:method}).

  \item \textbf{Về kiểm chứng thực nghiệm}: Thí nghiệm so sánh tám phương án khởi tạo
        ở cùng ngân sách (Mục~\ref{sec:exp-bakeoff}) xác nhận từng lựa chọn thiết kế:
        nhóm khởi tạo SALT vượt mọi đối chứng ngẫu nhiên ở cả mô hình hóa ngôn ngữ lẫn
        tác vụ hạ nguồn (F1 khoảng 66--68 so với 56 của đối chứng tốt nhất); chiếu
        theo từng token tốt hơn ánh xạ toàn cục; hệ số co $\gamma=0{,}1$ là cần thiết
        -- không co hoặc co không đủ khiến huấn luyện mắc kẹt và kết quả hạ nguồn tụt
        dưới cả đối chứng ngẫu nhiên; và căn chỉnh đóng băng cải thiện nhất quán.

  \item \textbf{Về hiệu quả của mô hình}: Sau 5 triệu tài liệu CulturaX, ViNeoBERT
        dẫn đầu 4/10 tác vụ trong bảng đánh giá đa nhiệm -- đọc hiểu UIT-ViQuAD
        (76,24 F1, hơn XLM-R-base 4,1 điểm), pseudo-perplexity, gán nhãn từ loại và
        XNLI -- và đạt kết quả bằng hoặc cao hơn XLM-R-base ở 8/10 tác vụ
        (Mục~\ref{sec:exp-mrc}). Đáng nói nhất, lượng tiếng Việt mô hình tiếp xúc chỉ
        khoảng 5 tỷ token -- ít hơn một đến hai bậc so với PhoBERT (40 epoch trên 20GB,
        bản v2 thêm 120GB OSCAR) hay XLM-R ($\sim$137GB tiếng Việt trong CC-100). Kết
        quả này xác nhận giả thuyết trung tâm của đề tài: \textit{tái sử dụng không
        gian biểu diễn sẵn có là con đường tiết kiệm để đưa các kiến trúc Encoder thế
        hệ mới đến với tiếng Việt}.
\end{itemize}

Bên cạnh mô hình, đề tài đóng góp một quy trình ba giai đoạn có thể tái lập -- khởi
tạo SALT cải tiến, căn chỉnh đóng băng, tiếp tục tiền huấn luyện theo lịch WSD với
khả năng huấn luyện theo phiên -- cùng bộ mã nguồn công khai đi kèm.

\section{Hạn chế của đề tài}
\label{sec:limitations}

Trong khuôn khổ thời gian và tài nguyên của một khóa luận, đề tài còn các hạn chế
sau:

\begin{itemize}
  \item \textbf{Ngân sách huấn luyện}: 5 triệu tài liệu vẫn nhỏ so với dữ liệu tiền
        huấn luyện của các mô hình đối sánh; ViNeoBERT chưa vượt PhoBERT-base-v2 ở
        nhóm tác vụ phân loại câu ngắn, dù xu hướng đường cong học cho thấy khoảng
        cách có thể tiếp tục thu hẹp khi tăng ngân sách.
  \item \textbf{Phạm vi đánh giá}: bảng đánh giá đa nhiệm hiện so sánh với hai mô
        hình cùng hạng \texttt{base}; các biến thể \texttt{large} và nhóm tác vụ truy
        hồi thông tin chưa được đưa vào so sánh đầy đủ.
  \item \textbf{Ngữ cảnh dài}: mô hình kế thừa trần ngữ cảnh 4.096 token của NeoBERT
        về mặt kiến trúc, nhưng quá trình tiếp tục tiền huấn luyện dùng chuỗi tối đa
        1.024 token và đề tài chưa kiểm chứng chất lượng ở ngữ cảnh dài.
  \item \textbf{Bộ tách từ}: từ vựng kế thừa từ ViDeBERTa hoạt động ở mức âm tiết;
        ảnh hưởng của việc không tách từ (so với PhoBERT dùng dữ liệu đã tách từ)
        chưa được phân tích riêng.
\end{itemize}

\section{Hướng phát triển}
\label{sec:future-work}

Từ các kết quả và hạn chế trên, chúng tôi đề xuất bốn hướng phát triển:

\begin{enumerate}
  \item \textbf{Mở rộng ngân sách tiếp tục tiền huấn luyện}: tăng số tài liệu CulturaX
        vượt mốc 5 triệu và khai thác thêm các pha làm nguội trên dữ liệu mới -- vốn
        đã cho thấy hiệu quả bật điểm ở mốc cuối -- nhằm vượt các mô hình họ PhoBERT
        trên các tác vụ hạ nguồn.
  \item \textbf{Kiểm chứng và khai thác ngữ cảnh dài}: huấn luyện bổ sung ở độ dài
        4.096 token và đánh giá trên các tác vụ đòi hỏi ngữ cảnh dài (hỏi đáp đa đoạn,
        truy hồi tài liệu); khi cần vượt trần hiện tại, các kỹ thuật nội suy RoPE như
        YaRN~\cite{peng2023yarn} là ứng viên tự nhiên.
  \item \textbf{Mở rộng bộ đánh giá}: bổ sung các tác vụ suy luận ngôn ngữ tự nhiên,
        gán nhãn chuỗi và truy hồi ngữ nghĩa cho tiếng Việt, nhằm đánh giá toàn diện
        hơn chất lượng biểu diễn của mô hình ở vai trò một Encoder đa dụng.
  \item \textbf{Tổng quát hóa quy trình}: áp dụng quy trình ba giai đoạn của đề tài
        cho các ngôn ngữ ít tài nguyên khác hoặc các kiến trúc nền khác, qua đó kiểm
        chứng tính tổng quát của các cải tiến (tập neo theo cặp dịch, khởi tạo lớp
        đầu ra tách rời, căn chỉnh đóng băng) ngoài phạm vi cặp Anh--Việt và NeoBERT.
\end{enumerate}
